% ==============================================================================
% SECTION 4: EVALUATION
% ==============================================================================
\section{Evaluation}
We evaluate ECAC through a series of experiments designed to validate its correctness properties and
assess performance under realistic conditions. Unlike many blockchain-based access control solutions that
focus mainly on throughput, our evaluation explores a broader range of metrics---from property
correctness to scalability, state growth, and revocation latency. All tests were conducted on a prototype
implementation of our event-sourced access control system deployed across multiple nodes. We used
synthetic event workloads and controlled network partitions to simulate challenging multi-issuer scenarios.

The evaluation is organized as follows:
\begin{itemize}
    \item \S\ref{subsec:methodology} describes our experimental methodology and workloads;
    \item \S\ref{subsec:property-validation} presents property validation results for the core guarantees;
    \item \S\ref{subsec:scalability} examines scalability and state growth (throughput, replay cost, storage);
    \item \S\ref{subsec:revocation-latency} measures revocation propagation delays in partitioned networks;
    \item \S\ref{subsec:limitations} candidly discusses observed weaknesses and limitations.
\end{itemize}

Throughout, we emphasize reproducibility by using deterministic workload generators with fixed seeds. All results are averaged over multiple runs to ensure statistical significance, and variance is reported where relevant.

\subsection{Experimental Methodology}\label{subsec:methodology}

\subsubsection{Implementation}
We implemented ECAC-core as a Rust library organized into four crates: \texttt{core} (operations, DAG, replay, policy, VCs, CRDTs), \texttt{store} (RocksDB persistence, checkpoints, audit), \texttt{net} (libp2p gossip, sync planning), and \texttt{cli} (command-line interface). The implementation uses Ed25519 for signatures, BLAKE3 for hashing (operations and state digests), XChaCha20-Poly1305 for authenticated encryption, and SHA-256 for status list bitstrings (per W3C VC specification). Operations are serialized in CBOR format for compactness and determinism. The prototype comprises approximately 8,500 lines of Rust code (excluding tests).

\subsubsection{Testbed Configuration}
Experiments were conducted on two testbed configurations:

\noindent\textbf{Local cluster.} A single machine with 16 cores (AMD Ryzen 9 5950X), 64\,GB RAM, and NVMe SSD storage, running Ubuntu 22.04 LTS with kernel 5.15. Multiple ECAC nodes execute as separate processes communicating over localhost. This configuration provides controlled, reproducible measurements with minimal network variance.

\noindent\textbf{Distributed cluster.} Ten virtual machines (8 vCPUs, 16\,GB RAM each) deployed across three availability zones in a cloud environment, connected via 10\,Gbps network with simulated latency (1--200\,ms) and packet loss (0--5\%) using \texttt{tc/netem}. This configuration validates behavior under realistic network conditions.

All experiments use Rust 1.85 in release mode with link-time optimization. Random seeds are fixed for reproducibility; we report the seed with each experiment.

\subsubsection{Workloads}
We designed five synthetic workloads to stress different aspects of ECAC:

\begin{table}[htbp]
\centering
\caption{Synthetic Workload Characteristics}
\label{tab:workloads}
\begin{tabular}{@{}lcccc@{}}
\toprule
\textbf{Workload} & \textbf{Writers} & \textbf{Revoke \%} & \textbf{Key Dist.} & \textbf{Purpose} \\
\midrule
W1: Linear      & 1  & 0\%  & Single   & Baseline \\
W2: Concurrent  & 10 & 0\%  & Single   & CRDT stress \\
W3: Partitioned & 10 & 0\%  & Disjoint & Sync efficiency \\
W4: Revocation  & 10 & 10\% & Uniform  & Policy overhead \\
W5: Mixed       & 20 & 2\%  & Zipfian  & Realistic \\
\bottomrule
\end{tabular}
\end{table}

\noindent\textbf{W1 (Linear chain).} A single authorized writer appends operations sequentially to a single key, establishing baseline throughput without concurrency or policy overhead.

\noindent\textbf{W2 (Concurrent writers).} Ten writers, each with valid credentials, concurrently modify the same key. This stresses the MVReg CRDT's conflict resolution and validates that concurrent updates converge deterministically.

\noindent\textbf{W3 (Partitioned keyspace).} Ten writers operate on disjoint key sets (no overlap). This measures sync efficiency when operations are causally independent and can be processed in parallel.

\noindent\textbf{W4 (Revocation-heavy).} Ten writers with 10\% of operations being Grant or Revoke events. Credentials are issued and revoked frequently, forcing continuous epoch recalculation during replay.

\noindent\textbf{W5 (Mixed realistic).} Twenty writers with Zipfian key distribution (skew factor 1.0), 2\% policy operations, and varying operation sizes. This approximates real-world collaborative editing patterns where some keys are hot spots.

Each workload scales from 10K to 10M operations. For multi-authority experiments, we configure 2--5 independent issuers, each managing credentials for a subset of writers.

\subsubsection{Metrics}
We measure the following primary metrics:

\begin{itemize}[leftmargin=1.5em]
    \item \textbf{Convergence rate}: Fraction of trials where all replicas produce identical state digests after processing the same operation set.
    \item \textbf{Replay throughput}: Operations processed per second during full replay ($\text{ops}/\text{s}$).
    \item \textbf{Replay latency}: Time to replay the complete log from genesis (ms).
    \item \textbf{Incremental replay speedup}: Ratio of full replay time to checkpoint-based incremental replay time.
    \item \textbf{Sync overhead}: Bytes transferred per operation during anti-entropy synchronization.
    \item \textbf{Storage growth}: Bytes per operation in the persistent store (ops + indexes + checkpoints).
    \item \textbf{Memory footprint}: Peak resident set size (RSS) during replay.
    \item \textbf{Policy evaluation time}: Microseconds to evaluate authorization for a single operation.
    \item \textbf{Revocation propagation delay}: Time from revocation issuance to enforcement at remote replicas.
    \item \textbf{Key rotation overhead}: Time and storage cost of per-tag encryption key rotation.
\end{itemize}

For latency metrics, we report p50, p95, and p99 percentiles. Statistical significance is established using Welch's t-test ($\alpha = 0.05$) with at least 30 trials per configuration.

\subsubsection{Experiment Design}
Our evaluation addresses five research questions mapped to the design properties P1--P8:

\begin{enumerate}[label=\textbf{RQ\arabic*},leftmargin=2.5em]
    \item \textit{Convergence under partitions} (P2): Do replicas that receive operations in different orders converge to identical states?
    \item \textit{Policy safety} (P3): Are unauthorized operations correctly rejected, and do revocations take effect deterministically?
    \item \textit{Scalability} (P8): How do replay time, storage, and memory scale with log size?
    \item \textit{Multi-authority conflicts}: When multiple issuers grant/revoke concurrently, does ECAC resolve conflicts deterministically and safely?
    \item \textit{Confidentiality overhead} (P7): What is the performance cost of per-tag encryption and key rotation?
\end{enumerate}

\subsection{Property Validation}\label{subsec:property-validation}

This section validates ECAC's core correctness guarantees through controlled experiments.

\subsubsection{Experiment E1: Convergence Under Random Ordering}
\textbf{Objective.} Verify that operation delivery order does not affect final state (Theorem 1).

\textbf{Method.} We generate a canonical operation log of 100K operations using workload W2 (10 concurrent writers). From this log, we create 100 distinct valid orderings by randomly permuting operations while respecting causal dependencies (parents must precede children). Each ordering is replayed independently on a fresh node. We compare final state digests (BLAKE3 hashes) across all replicas.

\textbf{Success criterion.} 100\% of orderings produce byte-identical state digests.

% ==============================================================================
% TODO: E1 RESULTS (Phase 1.3) - MUST HAVE
% ==============================================================================
\TODO{Insert E1 results after running experiment:}

\PLACEHOLDER{%
\textbf{Results.} All 100 orderings produced identical state digests (hash: \texttt{[INSERT\_HASH]}). Convergence rate: 100\% (100/100 trials). Mean replay time: [X] ms ($\sigma$ = [Y] ms). This confirms Theorem 1: the replay function is deterministic regardless of delivery order.
}

\subsubsection{Experiment E2: Cross-Platform Convergence}
\textbf{Objective.} Verify determinism across different hardware architectures and operating systems.

\textbf{Method.} We generate a canonical log of 1M operations and replay it on five platforms: Linux x86\_64 (Ubuntu 22.04), Linux ARM64 (Raspberry Pi 4, Ubuntu 22.04), macOS x86\_64 (Ventura), macOS ARM64 (M2), and Windows x86\_64 (Windows 11). All use Rust 1.85 with identical compiler flags.

\textbf{Success criterion.} All platforms produce identical 32-byte state digests.

% ==============================================================================
% TODO: E2 RESULTS (Phase 1.3) - SHOULD HAVE
% ==============================================================================
\TODO{Insert E2 results (lower priority than E1, E3, E6, E11):}

\PLACEHOLDER{%
\textbf{Results.} All five platforms produced identical state digests: \texttt{[INSERT\_32\_BYTE\_HASH]}. Replay times varied by platform (Linux x86\_64: [X] ms, ARM64: [Y] ms, macOS M2: [Z] ms), but final state was byte-identical across all architectures.
}

\subsubsection{Experiment E3: Revocation Correctness (Deny-Wins)}
\textbf{Objective.} Verify that revoked credentials cause subsequent operations to be skipped (Theorem 3).

\textbf{Method.}
\begin{enumerate}
    \item Issue a credential granting user $U$ write access to resource $R$ with scope $S$.
    \item User $U$ performs 1000 write operations to $R$.
    \item At operation 500, revoke $U$'s credential.
    \item Perform full replay from genesis.
    \item Verify: (a) operations 1--499 are applied, (b) operations 500--999 are skipped, (c) audit log records correct decisions.
\end{enumerate}

\textbf{Success criterion.} Zero operations from revoked epochs appear in final state; audit log shows 500 ``applied'' and 500 ``skipped (revoked)'' entries.

% ==============================================================================
% TODO: E3 RESULTS (Phase 1.3) - MUST HAVE
% ==============================================================================
\TODO{Insert E3 results after running experiment:}

\PLACEHOLDER{%
\textbf{Results.} The deny-wins mechanism correctly enforced revocation:
\begin{itemize}
    \item Operations 1--499: 499 applied (100\%)
    \item Operations 500--999: 500 skipped (100\%)
    \item Audit log entries: 499 ``applied'', 500 ``skipped (revoked)'', 1 ``grant'', 1 ``revoke''
    \item Final state contains exactly 499 writes from user $U$
\end{itemize}
This confirms Theorem 3: revoked credentials have no effect on final state.
}

\subsubsection{Experiment E4: Multi-Authority Conflict Resolution}
\textbf{Objective.} Verify deterministic resolution when multiple issuers make conflicting grants/revokes.

\textbf{Method.} Configure three independent issuers ($I_1$, $I_2$, $I_3$), each authorized to manage credentials for overlapping user sets. Simulate the following conflict scenario:
\begin{itemize}
    \item $I_1$ grants user $U$ role ``editor'' at $t_1$.
    \item $I_2$ grants user $U$ role ``editor'' at $t_2 > t_1$ (concurrent, different issuer).
    \item $I_3$ revokes user $U$'s role ``editor'' at $t_3 > t_2$.
    \item User $U$ writes at $t_4 > t_3$.
\end{itemize}
We inject these operations into 10 replicas in random orders and verify convergence.

\textbf{Success criterion.} All replicas: (a) accept $U$'s writes during $[t_1, t_3)$, (b) reject $U$'s write at $t_4$, (c) produce identical final digests.

% ==============================================================================
% TODO: E4 RESULTS (Phase 1.3) - SHOULD HAVE
% ==============================================================================
\TODO{Insert E4 results:}

\PLACEHOLDER{%
\textbf{Results.} All 10 replicas converged to identical state despite receiving operations in different orders:
\begin{itemize}
    \item Writes during $[t_1, t_3)$: accepted on all replicas
    \item Write at $t_4$: rejected on all replicas (skipped due to revocation)
    \item Final state digest: identical across all 10 replicas
    \item Convergence time after final sync: [X] ms (p50), [Y] ms (p95)
\end{itemize}
Multi-authority conflicts are resolved deterministically via the deny-wins rule.
}

\subsubsection{Experiment E5: Audit Integrity}
\textbf{Objective.} Verify tamper-evidence of the audit trail (Theorem 4).

\textbf{Method.} Generate a log of 50K operations, producing a hash-chained audit trail. Attempt three tampering scenarios: (a) delete an audit entry, (b) modify an entry's decision field, (c) reorder two entries. For each, verify that chain verification fails.

\textbf{Success criterion.} All tampering attempts detected with $>$ 99.999\% probability (collision resistance of BLAKE3).

% ==============================================================================
% TODO: E5 RESULTS (Phase 1.3) - SHOULD HAVE
% ==============================================================================
\TODO{Insert E5 results:}

\PLACEHOLDER{%
\textbf{Results.} All three tampering scenarios were detected:
\begin{itemize}
    \item Deletion: Chain verification failed at entry [N+1] (hash mismatch)
    \item Modification: Chain verification failed at modified entry (hash mismatch)
    \item Reordering: Chain verification failed at first reordered entry (hash mismatch)
\end{itemize}
Detection rate: 100\% across 1000 trials per scenario. This confirms Theorem 4.
}

\subsection{Scalability and Performance}\label{subsec:scalability}

\subsubsection{Experiment E6: Replay Time Scaling}
\textbf{Objective.} Characterize replay performance as log size grows (hypothesis: $O(n)$ in operation count).

\textbf{Method.} Generate logs of sizes 10K, 50K, 100K, 500K, 1M, 5M, and 10M operations using workload W5. Measure full replay time (cold start) and incremental replay time (from checkpoint at 90\% of log). Each configuration runs 30 trials.

\textbf{Analysis.} Fit linear regression $t = a \cdot n + b$; report $R^2$ and identify non-linear regimes if any.

% ==============================================================================
% TODO: E6 RESULTS (Phase 1.3) - MUST HAVE
% ==============================================================================
\TODO{Insert E6 results with table/figure:}

\PLACEHOLDER{%
\textbf{Results.} Replay time scales linearly with log size:

\begin{table}[h]
\centering
\caption{Replay Time vs. Log Size (Workload W5)}
\begin{tabular}{@{}rrrr@{}}
\toprule
\textbf{Operations} & \textbf{Full Replay (ms)} & \textbf{Incremental (ms)} & \textbf{Speedup} \\
\midrule
10K   & [X]   & [Y]   & [Z]$\times$ \\
50K   & [X]   & [Y]   & [Z]$\times$ \\
100K  & [X]   & [Y]   & [Z]$\times$ \\
500K  & [X]   & [Y]   & [Z]$\times$ \\
1M    & [X]   & [Y]   & [Z]$\times$ \\
5M    & [X]   & [Y]   & [Z]$\times$ \\
10M   & [X]   & [Y]   & [Z]$\times$ \\
\bottomrule
\end{tabular}
\end{table}

Linear regression: $t = [A] \cdot n + [B]$, $R^2 = [0.99+]$. Throughput: approximately [X] ops/s sustained.
}

\subsubsection{Experiment E7: Throughput Under Concurrency}
\textbf{Objective.} Measure sustainable operation append throughput with concurrent writers.

\textbf{Method.} Using workload W2 (10 concurrent writers) and W3 (10 partitioned writers), measure operations appended per second over a 60-second window. Vary writer count from 1 to 50.

\textbf{Metrics.} Peak throughput (ops/s), throughput at saturation, CPU utilization.

% ==============================================================================
% TODO: E7 RESULTS (Phase 1.3) - NICE TO HAVE
% ==============================================================================
\TODO{Insert E7 results:}

\PLACEHOLDER{%
\textbf{Results.} Peak throughput achieved with [N] concurrent writers:
\begin{itemize}
    \item W2 (single key): [X] ops/s peak, [Y] ops/s at saturation
    \item W3 (disjoint keys): [X] ops/s peak, [Y] ops/s at saturation
    \item CPU utilization at saturation: [Z]\%
\end{itemize}
Partitioned keyspace (W3) achieves [N]$\times$ higher throughput due to reduced CRDT conflict resolution overhead.
}

\subsubsection{Experiment E8: Storage Growth}
\textbf{Objective.} Characterize storage overhead per operation.

\textbf{Method.} Measure RocksDB storage size (all column families) after inserting 10K, 100K, 1M, and 10M operations. Break down by: (a) operations (CBOR bytes), (b) DAG edges, (c) verified VCs, (d) checkpoints, (e) audit log.

\textbf{Metrics.} Bytes per operation (total and per category); checkpoint compression ratio.

% ==============================================================================
% TODO: E8 RESULTS (Phase 1.3) - NICE TO HAVE
% ==============================================================================
\TODO{Insert E8 results:}

\PLACEHOLDER{%
\textbf{Results.} Storage grows linearly at approximately [X] bytes/operation:
\begin{itemize}
    \item Operations (CBOR): [X] bytes/op (average)
    \item DAG edges: [Y] bytes/op
    \item VCs + status lists: [Z] bytes/op (amortized)
    \item Checkpoints: [W] bytes total ([V]\% compression)
    \item Audit log: [U] bytes/op
\end{itemize}
At 10M operations: total storage = [X] GB.
}

\subsubsection{Experiment E9: Memory Profiling}
\textbf{Objective.} Characterize memory usage during replay.

\textbf{Method.} Profile peak RSS during full replay of logs from 10K to 10M operations. Identify major memory consumers (DAG structure, CRDT state, epoch index, pending queue).

\textbf{Metrics.} Peak RSS, memory per operation, growth rate.

% ==============================================================================
% TODO: E9 RESULTS (Phase 1.3) - NICE TO HAVE
% ==============================================================================
\TODO{Insert E9 results:}

\PLACEHOLDER{%
\textbf{Results.} Peak memory usage scales sub-linearly due to checkpointing:
\begin{itemize}
    \item 10K ops: [X] MB peak RSS
    \item 1M ops: [Y] MB peak RSS
    \item 10M ops: [Z] MB peak RSS
\end{itemize}
Major consumers: DAG structure ([A]\%), CRDT state ([B]\%), epoch index ([C]\%).
}

\subsubsection{Experiment E10: Checkpoint Efficiency}
\textbf{Objective.} Quantify the benefit of incremental replay from checkpoints.

\textbf{Method.} For a 1M operation log, create checkpoints at 25\%, 50\%, 75\%, and 90\% of log size. Measure incremental replay time to process remaining operations versus full replay.

\textbf{Metrics.} Speedup factor, checkpoint creation time, checkpoint size.

% ==============================================================================
% TODO: E10 RESULTS (Phase 1.3) - NICE TO HAVE
% ==============================================================================
\TODO{Insert E10 results:}

\PLACEHOLDER{%
\textbf{Results.} Checkpoints provide significant speedup:

\begin{table}[h]
\centering
\caption{Checkpoint Efficiency (1M operations)}
\begin{tabular}{@{}rrrr@{}}
\toprule
\textbf{Checkpoint} & \textbf{Incremental (ms)} & \textbf{Speedup} & \textbf{Size (MB)} \\
\midrule
25\%  & [X] & [Y]$\times$ & [Z] \\
50\%  & [X] & [Y]$\times$ & [Z] \\
75\%  & [X] & [Y]$\times$ & [Z] \\
90\%  & [X] & [Y]$\times$ & [Z] \\
\bottomrule
\end{tabular}
\end{table}

Checkpoint creation time: [X] ms. Compression ratio: [Y]:1.
}

\subsection{Revocation Propagation}\label{subsec:revocation-latency}

\subsubsection{Experiment E11: Revocation Under Partition}
\textbf{Objective.} Measure the window of vulnerability when revocation occurs during network partition.

\textbf{Method.} Deploy 6 nodes in two groups (A: 3 nodes, B: 3 nodes). Partition the network. Issue operations:
\begin{enumerate}
    \item Node in group A issues revocation for user $U$ at $t_r$.
    \item Disconnected node in group B (unaware of revocation) accepts $U$'s operations during partition.
    \item Heal partition; measure time until group B rejects $U$'s post-$t_r$ operations.
\end{enumerate}
Vary partition duration from 1 minute to 1 hour.

\textbf{Metrics.} Time to enforcement post-healing; number of unauthorized operations accepted during partition (subsequently invalidated by replay).

% ==============================================================================
% TODO: E11 RESULTS (Phase 1.3) - MUST HAVE
% ==============================================================================
\TODO{Insert E11 results:}

\PLACEHOLDER{%
\textbf{Results.} Revocation enforcement occurs within [X] seconds of partition healing:

\begin{table}[h]
\centering
\caption{Revocation Propagation Under Partition}
\begin{tabular}{@{}rrrr@{}}
\toprule
\textbf{Partition Duration} & \textbf{Ops During Partition} & \textbf{Time to Enforce (ms)} & \textbf{Ops Invalidated} \\
\midrule
1 min   & [X] & [Y] & [Z] \\
5 min   & [X] & [Y] & [Z] \\
15 min  & [X] & [Y] & [Z] \\
1 hour  & [X] & [Y] & [Z] \\
\bottomrule
\end{tabular}
\end{table}

Key finding: All operations accepted during partition were correctly invalidated upon replay after sync. The ``window of vulnerability'' equals partition duration plus sync time ([X] ms average).
}

\subsubsection{Experiment E12: Sync Convergence Time}
\textbf{Objective.} Measure time for partitioned replicas to converge after reconnection.

\textbf{Method.} Partition a 10-node cluster into two groups for 1 hour. Each group generates 10K operations independently. Heal partition and measure:
\begin{itemize}
    \item Time until all nodes report identical state digests.
    \item Total bytes transferred during sync.
    \item Operations per second during catch-up replay.
\end{itemize}

\textbf{Metrics.} Convergence time (p50, p95), sync bandwidth efficiency.

% ==============================================================================
% TODO: E12 RESULTS (Phase 1.3) - NICE TO HAVE
% ==============================================================================
\TODO{Insert E12 results:}

\PLACEHOLDER{%
\textbf{Results.} Convergence after 1-hour partition with 20K total operations:
\begin{itemize}
    \item Convergence time: [X] ms (p50), [Y] ms (p95)
    \item Bytes transferred: [Z] MB total ([W] bytes/op)
    \item Catch-up replay throughput: [V] ops/s
\end{itemize}
All 10 nodes achieved identical state digests within [X] seconds of partition healing.
}

\subsection{Confidentiality Overhead}\label{subsec:confidentiality}

\subsubsection{Experiment E13: Per-Tag Encryption Cost}
\textbf{Objective.} Measure the overhead of payload encryption for confidentiality.

\textbf{Method.} Compare throughput for workload W1 with and without per-tag encryption enabled. Measure:
\begin{itemize}
    \item Operation append latency (encryption time).
    \item Replay throughput (decryption time).
    \item Storage overhead (ciphertext expansion).
\end{itemize}

\textbf{Metrics.} Throughput reduction (\%), latency increase (ms), storage overhead (\%).

% ==============================================================================
% TODO: E13 RESULTS (Phase 1.3) - SHOULD HAVE
% ==============================================================================
\TODO{Insert E13 results:}

\PLACEHOLDER{%
\textbf{Results.} Encryption adds modest overhead:
\begin{itemize}
    \item Append throughput: [X] ops/s (encrypted) vs [Y] ops/s (plaintext) = [Z]\% reduction
    \item Replay throughput: [X] ops/s (encrypted) vs [Y] ops/s (plaintext) = [Z]\% reduction
    \item Storage overhead: [W]\% (XChaCha20-Poly1305 adds 40 bytes per payload)
    \item Encryption latency: [V] $\mu$s/op (p50)
\end{itemize}
}

\subsubsection{Experiment E14: Key Rotation Overhead}
\textbf{Objective.} Measure the cost of forward secrecy through key rotation.

\textbf{Method.} Simulate key rotation every 1000 operations (aggressive) and every 100K operations (conservative). For each rotation:
\begin{itemize}
    \item Measure KeyRotate operation creation time.
    \item Measure KeyGrant distribution time (re-encrypting new key to authorized users).
    \item Track storage growth from versioned keys.
\end{itemize}

\textbf{Metrics.} Rotation latency, key distribution overhead per authorized user, storage per rotation event.

% ==============================================================================
% TODO: E14 RESULTS (Phase 1.3) - NICE TO HAVE
% ==============================================================================
\TODO{Insert E14 results:}

\PLACEHOLDER{%
\textbf{Results.} Key rotation cost scales with number of authorized users:
\begin{itemize}
    \item KeyRotate creation: [X] ms
    \item KeyGrant per user: [Y] ms (public-key encryption)
    \item Storage per rotation: [Z] bytes + [W] bytes/user
\end{itemize}
At 1000-op rotation interval with 100 users: [V]\% throughput overhead. At 100K-op interval: $<$[U]\% overhead.
}

\subsection{Limitations and Threats to Validity}\label{subsec:limitations}

\subsubsection{Inherent Limitations}
ECAC's design involves fundamental trade-offs:

\noindent\textbf{No immediate revocation for offline nodes.} By design, a disconnected node cannot learn of revocations until it syncs. During partition, a node may accept operations from revoked users; these are invalidated upon replay after reconnection, but the window exists. This is an unavoidable consequence of offline-first operation without consensus.

\noindent\textbf{Append-only storage growth.} The event log grows monotonically. While checkpointing reduces replay cost, the full log must be retained for auditability. Long-running deployments require storage planning or log compaction with trusted witnesses.

\noindent\textbf{Replay cost after long partitions.} Nodes that have been offline for extended periods must replay potentially large operation sets upon reconnection. Checkpoints mitigate but do not eliminate this cost.

\noindent\textbf{Confidentiality vs.\ full state access.} The current implementation supports per-tag encryption, but metadata (operation structure, timestamps, author public keys) remains visible to all replicas. Full metadata privacy would require additional cryptographic machinery (e.g., encrypted indexes).

\subsubsection{Threats to Validity}

\noindent\textbf{Internal validity.} Implementation bugs may affect results. We mitigate this through extensive unit testing (1,200+ tests), property-based testing with \texttt{proptest}, and code review. Measurement overhead is minimized by separating timing instrumentation from hot paths.

\noindent\textbf{External validity.} Synthetic workloads may not reflect production patterns. We include workload W5 (Zipfian distribution) to approximate real-world access patterns, but validation on production traces remains future work.

\noindent\textbf{Construct validity.} Our metrics may not capture all relevant properties. We focus on correctness (convergence, policy safety) and efficiency (throughput, latency, storage) as primary concerns for access control systems.

\noindent\textbf{Reproducibility.} All experiments use fixed random seeds, pinned toolchain versions, and automated scripts. We provide a reproducibility package with build instructions and expected outputs.
